\documentclass[10pt,a4paper]{amsart}
\usepackage[margin=19mm]{geometry}
\usepackage{amsmath,amssymb,mathtools}
\usepackage[scr=rsfs,cal=euler]{mathalfa}
\usepackage{xfrac}
\usepackage[unicode,hidelinks,bookmarks=false]{hyperref}
\newcommand{\ie}{i.e.\ }
\newcommand{\cf}{cf.\ }
\newcommand{\resp}{respectively\ }
\newcommand{\Subsection}{\subsection}
\providecommand{\nobreakdash}{\nobreak\hbox{-}}
\setlength{\emergencystretch}{2em}
\multlinegap0pt
\usepackage{bm}
\usepackage{enumerate}
\usepackage{amssymb}
\usepackage{dsfont}
\newtheorem{proposition}{Proposition}
\title{Global Gevrey Hypoellipticity of Involutive Systems on Non-Compact Manifolds}
\author{Sandro Coriasco, Alexandre Kirilov, Wagner Augusto Almeida de Moraes, Pedro Meyer Tokoro}
\date{Source: arXiv:2602.16366v1 (CC BY 4.0)}
\begin{document}
\maketitle

\noindent\emph{Source and licence.} Global Gevrey Hypoellipticity of Involutive Systems on Non-Compact Manifolds. Version arXiv:2602.16366v1 (CC BY 4.0), distributed by arXiv under the Creative Commons Attribution 4.0 International licence, http://creativecommons.org/licenses/by/4.0/, as recorded in the arXiv licence metadata for this exact version and shown on the abstract page https://arxiv.org/abs/2602.16366v1. Full licence notice: \texttt{licenses/arxiv-2602.16366-CC-BY-4.0.txt}.\par\smallskip

\noindent\textit{Editorial scope.} Counted unit: the article's proposition liou\_DCs (source0.tex lines 691--701). The article's complete original proof of it is delivered with it (source0.tex lines 703--1027, one \texttt{proof} environment of 325 lines). No mathematics is written, repaired or re-derived: the statement and the proof are the article's own lines, in the article's own notation, and no retained line is substituted or re-flowed.  Retained uncounted as the source-local context the proof needs: lines 468--470; lines 479--485; lines 603--609; lines 615--621.  Nothing was omitted from any retained span, so the package carries no omission record.  No retained line contains a citation, so the wrapper carries no bibliography and no bibliography entry.  No retained line contains a figure environment, an image-inclusion command or drawing code, so no image is delivered and the package declares no asset dependency.  Editorial formatting only, in addition: an AMS \texttt{amsart} preamble that restates the article's own declarations and macros used by the retained lines, namely the environments \texttt{proposition} on separate counters, together with the standard packages those lines need.  The retained environments print in this document as Proposition 1 in order of appearance, whereas the article prints the same items with the numbering of its own full document.  The complete native source of the article as distributed in the arXiv e-print for this exact version is bundled unchanged as \texttt{sources/194876/source0.tex}.
	\begin{proposition}\label{gh_lapl}
		Fix $s>1$ and let $g$ be a ${[s]}$-scattering metric on the interior $M$ of a compact manifold with boundary $\overline{M}$. Then, the Laplace--Beltrami operator $\Delta_g$ is ${[s]}$-globally hypoelliptic on $M$.
	\end{proposition}

	\begin{proposition}\label{gev_gh}
		Fix $s>1$. Let $M$ be the interior of a compact manifold with boundary, endowed with a ${[s]}$-scattering metric $g$. Then the exterior derivative $\mathrm{d}:\mathscr{G}^{[s]}(M)\to\mathsf{\Lambda}^1\mathscr{G}^{[s]}(M)$ is ${[s]}$-globally hypoelliptic. Furthermore, for any $\xi \in \mathbb{Z}^m$, the operator $\mathbb{L}_\xi:\mathscr{G}^{[s]}(M)\to\mathsf{\Lambda}^1\mathscr{G}^{[s]}(M)$ defined by
		\[
		\mathbb{L}_\xi u = \mathrm{d}u + i(\xi\cdot\boldsymbol{\omega})\wedge u
		\]
		is also ${[s]}$-globally hypoelliptic.
	\end{proposition}

			\begin{equation}\label{DCs}
				|\eta+\boldsymbol{A}\xi|
				\geq
				C_\varepsilon
				e^{-\varepsilon (|\eta|+|\xi|)^{1/s}},
				\tag{DC$_s$}
			\end{equation}

			\begin{equation}\label{DCss}
				|\eta+\boldsymbol{A}\xi|
				\geq
				C
				e^{-\varepsilon (|\eta|+|\xi|)^{1/s}},
				\tag{DC$_{(s)}$}
			\end{equation}

	\begin{proposition}\label{liou_DCs}
		Let $\boldsymbol{\omega}=(\omega_1,\dots,\omega_m)$ be a family of $1$-forms in 
		$\mathsf{\Lambda}^1\mathscr{G}^{[s]}_{\partial M}(M)$. Then, the system $\boldsymbol{\omega}$ is:
		\begin{enumerate}
			\item[(i)] $s$-exponential Liouville if and only if it is not rational and 
			$A(\boldsymbol{\omega})$ does not satisfy \eqref{DCs};
			
			\item[(ii)] $(s)$-exponential Liouville if and only if it is not rational and 
			$A(\boldsymbol{\omega})$ does not satisfy \eqref{DCss}.
		\end{enumerate}
	\end{proposition}

	\begin{proof}
	By the Hodge theorem for scattering manifolds, there exist 
	$\vartheta_1,\dots,\vartheta_d \in \mathcal{H}^1(M)$ whose cohomology classes form a basis of 
	$H^1_{\partial M}(M)$ dual to the cycles $\sigma_1,\dots,\sigma_d$. 
	Accordingly, for each $k=1,\dots,m$, we can write
	\begin{equation}\label{decomp_omega_k}
		\omega_k = \sum_{\ell=1}^d \lambda_{\ell k}\,\vartheta_\ell
		+ \mathrm{d}v_k,
	\end{equation}
	for suitable functions $v_k$.  By Propositions~\ref{gev_gh} and~\ref{gh_lapl}, we have $v_k\in\mathscr{G}^{[s]}(M)$ and $\vartheta_\ell \in\mathsf{\Lambda}^1 \mathscr{G}^{[s]}(M)$.
	
	Integrating \eqref{decomp_omega_k} over $\sigma_\ell$ yields
	\[
	\lambda_{\ell k} = A(\boldsymbol{\omega})_{\ell k}, \qquad \ell=1,\dots,d,\ k=1,\dots,m.
	\]
	
	Consequently, for every $\xi\in\mathbb{Z}^m$,
	\begin{equation}\label{omega}
		\xi\cdot\boldsymbol{\omega} = \sum_{\ell=1}^d \left( \sum_{k=1}^m A(\boldsymbol{\omega})_{\ell k}\xi_k \right)\vartheta_\ell + \sum_{k=1}^m \xi_k\,\mathrm{d}v_k.
	\end{equation}
	
	\medskip
	
	\noindent
	$(i)$ Suppose that $A(\boldsymbol{\omega})$ does not satisfy \eqref{DCs}. 
	Then there exist $\varepsilon>0$ and a sequence
	\(\{(\eta^{(j)},\xi^{(j)})\}\) in \(\mathbb{Z}^d \times (\mathbb{Z}^m \setminus \{0\})\) such that
	\begin{equation}\label{not_DC_r}
		|\eta^{(j)} + A(\boldsymbol{\omega})\xi^{(j)}| < j^{-1} e^{-\varepsilon |\xi^{(j)}|^{1/s}}, \qquad j\in\mathbb{N}.
	\end{equation}
	
	We claim that $|\xi^{(j)}|\to\infty$. 
	Indeed, if $\{\xi^{(j)}\}_{j\in\mathbb{N}}$ is bounded, then $A(\boldsymbol{\omega})\xi^{(j)}$ would take only finitely many values. 
	Since
	\[
	|\eta^{(j)}| \le |\eta^{(j)} + A(\boldsymbol{\omega})\xi^{(j)}| + |A(\boldsymbol{\omega})\xi^{(j)}|,
	\]
	the sequence $\{\eta^{(j)}\}_{j\in\mathbb{N}}$ would also be bounded, and hence $(\eta^{(j)},\xi^{(j)})$ would assume only finitely many values, contradicting \eqref{not_DC_r}. 
	Thus $|\xi^{(j)}|\to\infty$.
	
	For each $j\in\mathbb{N}$, define
	\begin{equation}\label{vartheta_r}
		\theta_j = -\sum_{\ell=1}^d \eta^{(j)}_\ell\,\vartheta_\ell + \sum_{k=1}^m \xi^{(j)}_k\,\mathrm{d}v_k \in		\mathsf{\Lambda}^1 \mathscr{G}^s(M),
	\end{equation}
	where
	\[\xi^{(j)}=(\xi^{(j)}_1,\dots,\xi^{(j)}_m)\quad\text{and}\quad \eta^{(j)}=(\xi^{(j)}_1,\dots,\eta^{(j)}_d).\]
	
	Each $\theta_j$ is integral, since $(\eta^{(j)})_\ell\in\mathbb{Z}$, 
	the forms $\vartheta_\ell$ are integral, and exact forms are integral.
	
	Using \eqref{omega} and \eqref{vartheta_r}, we obtain
	\begin{align*}
		\rho_j & := e^{\varepsilon|\xi^{(j)}|^{\frac{1}{s}}}(\xi^{(j)} \cdot \boldsymbol{\omega} - \theta_j)\\
		& = e^{\varepsilon|\xi^{(j)}|^{\frac{1}{s}}}\left[\sum_{\ell=1}^{d}\left(\sum_{k=1}^{m}A(\boldsymbol{\omega})_{\ell k}\xi^{(j)}_k\right)\vartheta_\ell + \sum_{k=1}^{m}\xi^{(j)}_k\,\mathrm{d}v_k + \sum_{\ell=1}^{d}\eta^{(j)}_\ell\vartheta_\ell - \sum_{k=1}^{m}\xi^{(j)}_k\,\mathrm{d}v_k\right]\\
		& = e^{\varepsilon|\xi^{(j)}|^{\frac{1}{s}}} \sum_{\ell=1}^{d} \left(\sum_{k=1}^{m} A(\boldsymbol{\omega})_{\ell k}\xi^{(j)}_k + \eta^{(j)}_\ell\right)\vartheta_\ell.
	\end{align*}
	
	
	It remains to prove that the sequence $\{\rho_j\}_{j\in\mathbb{N}}$ is bounded in 	$\mathsf{\Lambda}^1 \mathscr{G}^s(M)$. Let $(V,t)$ be a local chart and let $K\Subset V$.  For each multi-index $\alpha\in\mathbb{N}_0^n$, we have 
	\begin{equation}\label{sup_dtrho_r}
		\sup_{t\in K} |\partial_t^\alpha \rho_j(t)| 
		= e^{\varepsilon|\xi^{(j)}|^{1/s}} \sum_{\ell=1}^d 	\left| \sum_{k=1}^m
		A(\boldsymbol{\omega})_{\ell k}\xi^{(j)}_k + \eta^{(j)}_\ell \right|
		\, 	\sup_{t\in K} |\partial_t^\alpha \vartheta_\ell(t)|.
	\end{equation}
	
	Since each $\vartheta_\ell$ is Gevrey of order $s$, 
	there exist constants $C,h>0$, depending only on $K$, such that
	\begin{equation}\label{est_harm_r}
		\sup_{t\in K}	|\partial_t^\alpha \vartheta_\ell(t)|
		\le C h^{|\alpha|}\alpha!^s, \qquad \alpha\in\mathbb{N}_0^n,\ \ell=1,\dots,d.
	\end{equation}
	
	Substituting \eqref{est_harm_r} into \eqref{sup_dtrho_r}, we obtain
	\begin{align*}
		\sup_{t\in K} |\partial_t^\alpha \rho_j(t)|
		&\le C h^{|\alpha|}\alpha!^s e^{\varepsilon|\xi^{(j)}|^{1/s}} \sum_{\ell=1}^d
		\left| \sum_{k=1}^m A(\boldsymbol{\omega})_{\ell k}\xi^{(j)}_k + \eta^{(j)}_\ell 	\right|.
	\end{align*}
	
	Observe that
	\[
	\sum_{\ell=1}^d \left| \sum_{k=1}^m A(\boldsymbol{\omega})_{\ell k}\xi^{(j)}_k +
	\eta^{(j)}_\ell \right|
	\le C_1 |\eta^{(j)} + A(\boldsymbol{\omega})\xi^{(j)}|, 
	\]
	for some constant $C_1>0$ depending only on $d$.  Hence,
	\[
	\sup_{t\in K} |\partial_t^\alpha \rho_j(t)| \le C' h^{|\alpha|}\alpha!^s 
	e^{\varepsilon|\xi^{(j)}|^{1/s}} |\eta^{(j)} + A(\boldsymbol{\omega})\xi^{(j)}|.
	\]
	
	By \eqref{not_DC_r}, the sequence
	\[
	\big\{ e^{\varepsilon|\xi^{(j)}|^{1/s}} |\eta^{(j)} + A(\boldsymbol{\omega})\xi^{(j)}|	\big\}_{j\in\mathbb{N}}
	\]
	is bounded in $\mathbb{R}$. Therefore there exists $C''>0$, independent of $j\in\mathbb{N}$, such that
	\[
	\sup_{t\in K} |\partial_t^\alpha \rho_j(t)| 
	\le C'' h^{|\alpha|}\alpha!^s, \qquad \alpha\in\mathbb{N}_0^n.
	\]
	
	This shows that $\{\rho_j\}_{j\in\mathbb{N}}$ is bounded in $\mathsf{\Lambda}^1\mathscr{G}^s(M)$.
	
	
	\medskip
	
	Conversely, suppose that the family $\boldsymbol{\omega}$ is $s$-exponential Liouville.  Then there exist $\varepsilon>0$, a sequence $\{\xi^{(j)}\}_{j\in\mathbb{N}} $ in $\mathbb{Z}^m$ with $|\xi^{(j)}|\to\infty$, and a sequence of integral $1$-forms $\{\theta_j\}_{j\in\mathbb{N}}$ such that the sequence
	\[
	\{\rho_j\}_{j\in\mathbb{N}} := \big\{e^{\varepsilon|\xi^{(j)}|^{1/s}} (\xi^{(j)} \cdot \boldsymbol{\omega} - \theta_j)\big\}_{j\in\mathbb{N}}
	\]
	is bounded in $\mathsf{\Lambda}^1 \mathscr{G}^s(M)$.
	
	Since each $\theta_j$ is integral and belongs to $\mathsf{\Lambda}^1\mathscr{G}^s(M)$, we may write
	\begin{equation}\label{theta_r}
		\theta_j = \sum_{\ell=1}^{d} \mu_{j\ell}\,\vartheta_\ell + \mathrm{d}w_j,
	\end{equation}
	with $\mu_{j\ell} \in \mathbb{Z}$ and $w_j \in \mathscr{G}^s(M)$.
	
	Combining \eqref{omega} and \eqref{theta_r}, we obtain
	\begin{align}
		\rho_j &= e^{\varepsilon|\xi^{(j)}|^{1/s}} \Bigg[ \sum_{\ell=1}^{d}
		\left(\sum_{k=1}^{m} A(\boldsymbol{\omega})_{\ell k}\xi^{(j)}_k
		- \mu_{j\ell} \right) \vartheta_\ell + \sum_{k=1}^{m}\xi^{(j)}_k\,\mathrm{d}v_k
		- \mathrm{d}w_j \Bigg]. \label{rho_j_r}
	\end{align}
	
	Since $\int_{\sigma_\ell} \mathrm{d}v_k = 0$ and $\int_{\sigma_\ell} \mathrm{d}w_j = 0$, 
	integration of \eqref{rho_j_r} over $\sigma_\ell$ yields
	\[
	\frac{1}{2\pi} \int_{\sigma_\ell} \rho_j = e^{\varepsilon|\xi^{(j)}|^{1/s}} \left( \sum_{k=1}^{m}
	A(\boldsymbol{\omega})_{\ell k}\xi^{(j)}_k - \mu_{j\ell} \right), \qquad \ell=1,\dots,d.
	\]

	Since $\{\rho_j\}_{j\in\mathbb{N}}$ is bounded in $\mathsf{\Lambda}^1 \mathscr{G}^s(M)$, there exists a constant $C>0$ such that
	\[
	\left| \int_{\sigma_\ell} \rho_j \right| \le C, \qquad \forall\, j\in\mathbb{N},\ \ell=1,\dots,d.
	\]
	
	Therefore,
	\[
	e^{\varepsilon|\xi^{(j)}|^{1/s}} \left| \sum_{k=1}^{m} A(\boldsymbol{\omega})_{\ell k}\xi^{(j)}_k - \mu_{j\ell} \right| \le C', \qquad \ell=1,\dots,d,
	\]
	for some $C'>0$ independent of $j$.
	
	Defining $\eta^{(j)} := -(\mu_{j1}, \dots, \mu_{jd}) \in \mathbb{Z}^d$, we obtain
	\[
	|\eta^{(j)} + A(\boldsymbol{\omega})\xi^{(j)}| \le C'' e^{-\varepsilon|\xi^{(j)}|^{1/s}},
	\]
	for some $C''>0$ independent of $j$.
	
	Since $|\xi^{(j)}| \to \infty$, passing to a subsequence if necessary, we can assume that $e^{-\frac{\varepsilon}{2}|\xi^{(j)}|^{\frac{1}{s}}} \leq (C''j)^{-1}$. Hence,
	\[
	|\eta^{(j)} + A(\boldsymbol{\omega}) \xi^{(j)}| \leq j^{-1} e^{-\frac{\varepsilon}{2}|\xi^{(j)}|^{\frac{1}{s}}}, \qquad \forall\, j \in \mathbb{N},
	\]
	which implies that \eqref{DCs} does not hold.
	
	\medskip
	
	\noindent
	$(ii)$ Suppose first that $A(\boldsymbol{\omega})$ does not satisfy \eqref{DCss}. 
	Then, for every $j\in\mathbb{N}$, there exist  	$(\eta^{(j)},\xi^{(j)})$ in $\mathbb{Z}^d \times (\mathbb{Z}^m\setminus\{0\})$ such that
	\begin{equation}\label{not_DC_b}
		|\eta^{(j)} + A(\boldsymbol{\omega})\xi^{(j)}| 	< e^{-j|\xi^{(j)}|^{1/s}}.
	\end{equation}
	
	As in the proof of (i) above, one verifies that $\{|\xi^{(j)}|\}$ is unbounded. Hence, passing to a subsequence if necessary, we may assume that $|\xi^{(j)}|\to\infty$.
	
	For each $j$, define
	\begin{equation*}
		\theta_j = -\sum_{\ell=1}^d \eta^{(j)}_\ell\,\vartheta_\ell + \sum_{k=1}^m \xi^{(j)}_k\,\mathrm{d}v_k,
	\end{equation*}
	where
	\[\xi^{(j)}=(\xi^{(j)}_1,\dots,\xi^{(j)}_m)\quad\text{and}\quad \eta^{(j)}=(\xi^{(j)}_1,\dots,\eta^{(j)}_d)\]
	as before.	
	Each $\theta_j$ is integral, since $\eta^{(j)}_\ell\in\mathbb{Z}$, $\ell=1,\dots,d$, and exact forms are integral.
	
	
	Fix $\varepsilon>0$. We compute
	\begin{align*}
		\rho_j^{(\varepsilon)} 
		&:= e^{\varepsilon|\xi^{(j)}|^{\frac{1}{s}}} (\xi^{(j)} \cdot \boldsymbol{\omega} - \theta_j) \\
		&=
		e^{\varepsilon|\xi^{(j)}|^{\frac{1}{s}}} \left[ \sum_{\ell=1}^{d} \left(
		\sum_{k=1}^{m} A(\boldsymbol{\omega})_{\ell k}\xi^{(j)}_k \right)\vartheta_\ell
		+ \sum_{k=1}^{m}\xi^{(j)}_k\,\mathrm{d}v_k + 	\sum_{\ell=1}^{d} \eta^{(j)}_\ell\vartheta_\ell - \sum_{k=1}^{m}\xi^{(j)}_k\,\mathrm{d}v_k \right]\\
		& = e^{\varepsilon|\xi^{(j)}|^{\frac{1}{s}}} \sum_{\ell=1}^{d}
		\left( 	\sum_{k=1}^{m} A(\boldsymbol{\omega})_{\ell k}\xi^{(j)}_k + \eta^{(j)}_\ell
		\right)\vartheta_\ell. 
	\end{align*}
	
	
	We now prove that the sequence $\{\rho_j^{(\varepsilon)}\}_{j\in\mathbb{N}}$ is bounded in 
	$\mathsf{\Lambda}^1\mathscr{G}^{(s)}(M)$.
	
	Let $(V,t)$ be a local chart and let $K\Subset V$. 
	Since each $\vartheta_\ell$ is Gevrey of Beurling type of order $s$, 
	for every $h>0$ there exists $C_h>0$ such that
	\[
	\sup_{t\in K} |\partial_t^\alpha \vartheta_\ell(t)|
	\le C_h\, h^{|\alpha|}\alpha!^s, \qquad \alpha\in\mathbb{N}_0^n,
	\ \ell=1,\dots,d.
	\]
	
	Therefore, for every $\alpha\in\mathbb{N}_0^n$, we have
	\begin{align*}
		\sup_{t\in K} |\partial_t^\alpha \rho_j^{(\varepsilon)}(t)|
		&\le e^{\varepsilon|\xi^{(j)}|^{\frac{1}{s}}} \sum_{\ell=1}^{d} \left|
		\sum_{k=1}^{m} A(\boldsymbol{\omega})_{\ell k}\xi^{(j)}_k  + \eta^{(j)}_\ell \right| \sup_{t\in K} |\partial_t^\alpha \vartheta_\ell(t)| \\
		&\le C_h\, h^{|\alpha|}\alpha!^s e^{\varepsilon|\xi^{(j)}|^{\frac{1}{s}}} 		\sum_{\ell=1}^{d} \left| \sum_{k=1}^{m} A(\boldsymbol{\omega})_{\ell k}\xi^{(j)}_k
		+ \eta^{(j)}_\ell \right|.
	\end{align*}
	
	Since
	\[
	\sum_{\ell=1}^{d} \left| \sum_{k=1}^{m} A(\boldsymbol{\omega})_{\ell k}\xi^{(j)}_k
	+ \eta^{(j)}_\ell \right| \le C_1\, |\eta^{(j)} + A(\boldsymbol{\omega})\xi^{(j)}| 
	\]
	for some constant $C_1>0$ depending only on $d$, 
	we obtain
	\[
	\sup_{t\in K} |\partial_t^\alpha \rho_j^{(\varepsilon)}(t)|  \le C_h'\, h^{|\alpha|}\alpha!^s e^{\varepsilon|\xi^{(j)}|^{\frac{1}{s}}} |\eta^{(j)} + A(\boldsymbol{\omega})\xi^{(j)}|.
	\]
	
	Using \eqref{not_DC_b}, we deduce
	\[
	\sup_{t\in K} |\partial_t^\alpha \rho_j^{(\varepsilon)}(t)|
	\le 	C_h'\, h^{|\alpha|}\alpha!^s e^{(\varepsilon-j)|\xi^{(j)}|^{\frac{1}{s}}}.
	\]
	
	If $j>\varepsilon$, then $e^{(\varepsilon-j)|\xi^{(j)}|^{\frac{1}{s}}}\le 1$, and hence
	\[
	\sup_{t\in K} |\partial_t^\alpha \rho_j^{(\varepsilon)}(t)|
	\le C_h''\, h^{|\alpha|}\alpha!^s,
	\]
	uniformly in $j$, for $j$ sufficiently large.
	
	Since $\varepsilon>0$ was arbitrary and the above estimate holds for every $h>0$, 
	it follows that $\{\rho_j^{(\varepsilon)}\}_{j\in\mathbb{N}}$ is bounded in 
	$\mathsf{\Lambda}^1\mathscr{G}^{(s)}(M)$.	
	Therefore, $\boldsymbol{\omega}$ is $(s)$-exponential Liouville.
	
	Conversely, suppose that $\boldsymbol{\omega}$ is $(s)$-exponential Liouville. 
	Then there exists a sequence $\{\xi^{(j)}\}_{j\in\mathbb{N}}$ in $\mathbb{Z}^m$ with $|\xi^{(j)}|\to\infty$ 
	and integral $1$-forms $\theta_j$ such that, for every $\varepsilon>0$, 
	the sequence
	\[
	\{\rho_j^{(\varepsilon)} \}_{j\in\mathbb{N}} = \big\{e^{\varepsilon|\xi^{(j)}|^{1/s}} 	(\xi^{(j)}\cdot \boldsymbol{\omega} - \theta_j)\big\}_{j\in\mathbb{N}}
	\]
	is bounded in $\mathsf{\Lambda}^1\mathscr{G}^{(s)}(M)$.
	
	Writing
	\[
	\theta_j = \sum_{\ell=1}^d \mu_{j\ell}\,\vartheta_\ell + \mathrm{d}w_j,
	\qquad \mu_{j\ell}\in\mathbb{Z},
	\]
	and integrating over $\sigma_\ell$, we obtain
	\[
	\frac{1}{2\pi} \int_{\sigma_\ell} \rho_j^{(\varepsilon)} 
	= e^{\varepsilon|\xi^{(j)}|^{1/s}} \left( A(\boldsymbol{\omega})\xi^{(j)} - \mu_j \right)_\ell.
	\]
	
	Since $\{\rho_j^{(\varepsilon)}\}_{j\in\mathbb{N}}$ is bounded in $\mathsf{\Lambda}^1 \mathscr{G}^{(s)}(M)$, for every $\varepsilon>0$ there exists $C_\varepsilon>0$ such that
	\[
	|\eta^{(j)} + A(\boldsymbol{\omega})\xi^{(j)}| 	\le C_\varepsilon e^{-\varepsilon|\xi^{(j)}|^{1/s}},
	\qquad \eta^{(j)}:=-\mu_j\in\mathbb{Z}^d.
	\]
	
	As this estimate holds for every $\varepsilon>0$,  condition \eqref{DCss} cannot be satisfied. 
	
	Conversely, suppose that the family $\boldsymbol{\omega}$ is $(s)$-exponential Liouville. 	Then there exist a sequence $\{\xi^{(j)}\}_{j\in\mathbb{N}}$ in $\mathbb{Z}^m$, with $|\xi^{(j)}|\to\infty$, and a sequence of integral $1$-forms $\{\theta_j\}_{j\in\mathbb{N}}$ such that, 
	for every $\varepsilon>0$, the sequence
	\[
	\{\rho_j^{(\varepsilon)}\}_{j\in\mathbb{N}} = \big\{e^{\varepsilon|\xi^{(j)}|^{\frac{1}{s}}} 	(\xi^{(j)} \cdot \boldsymbol{\omega} - \theta_j)\big\}_{j\in\mathbb{N}}
	\]
	is bounded in $\mathsf{\Lambda}^1 \mathscr{G}^{(s)}(M)$.
	
	Since each $\theta_j$ is integral and belongs to 
	$\mathsf{\Lambda}^1\mathscr{G}^{(s)}(M)$, 
	we may write
	\begin{equation}\label{theta_bss}
		\theta_j = \sum_{\ell=1}^{d} \mu_{j\ell}\,\vartheta_\ell + \mathrm{d}w_j,
	\end{equation}
	where $\mu_{j\ell}\in\mathbb{Z}$ and $w_j\in\mathscr{G}^{(s)}(M)$.
	
	Fix $\varepsilon>0$. 
	Combining \eqref{omega} and \eqref{theta_bss}, we obtain
	\begin{align*}
		\rho_j^{(\varepsilon)}
		&= e^{\varepsilon|\xi^{(j)}|^{\frac{1}{s}}} \left[ 	\sum_{\ell=1}^{d} \left( \sum_{k=1}^{m} A(\boldsymbol{\omega})_{\ell k}\xi^{(j)}_k \right)\vartheta_\ell
		+ \sum_{k=1}^{m}\xi^{(j)}_k\,\mathrm{d}v_k 	- \sum_{\ell=1}^{d} \mu_{j\ell}\,\vartheta_\ell - \mathrm{d}w_j
		\right] \\ 
		&= e^{\varepsilon|\xi^{(j)}|^{\frac{1}{s}}} \sum_{\ell=1}^{d} \left( \sum_{k=1}^{m}
		A(\boldsymbol{\omega})_{\ell k}\xi^{(j)}_k - \mu_{j\ell} 	\right)\vartheta_\ell
		+ e^{\varepsilon|\xi^{(j)}|^{\frac{1}{s}}} 	\left( 	\sum_{k=1}^{m} \xi^{(j)}_k\,\mathrm{d}v_k - \mathrm{d}w_j \right).
	\end{align*}
	
	Integrating over the cycles $\sigma_\ell$,
	we obtain
	\[
	\frac{1}{2\pi} \int_{\sigma_\ell} \rho_j^{(\varepsilon)} =  	e^{\varepsilon |\xi^{(j)}|^{\frac{1}{s}}}\left( \sum_{k=1}^{m} A(\boldsymbol{\omega})_{\ell k}\xi^{(j)}_k
	- \mu_{j\ell} 	\right), 	\qquad 	\ell=1,\dots,d.
	\]
	
	Since $\{\rho_j^{(\varepsilon)}\}_{j\in\mathbb{N}}$ is bounded in  $\mathsf{\Lambda}^1 \mathscr{G}^{(s)}(M)$, it is bounded in each seminorm defining the Beurling topology. In particular, there exists a constant 
	$C_\varepsilon>0$ such that
	\[
	\left| 	\int_{\sigma_\ell} 	\rho_j^{(\varepsilon)} 	\right|
	\le 	C_\varepsilon, 	\qquad 	\forall\, j\in\mathbb{N},\  \ell=1,\dots,n.
	\]
	
	Therefore,
	\[
	e^{\varepsilon|\xi^{(j)}|^{\frac{1}{s}}} \left| \sum_{k=1}^{m} 	A(\boldsymbol{\omega})_{\ell k}\xi^{(j)}_k - \mu_{j\ell} 	\right|
	\le C_\varepsilon, 	\qquad	\forall\,j\in\mathbb{N},\ \ell=1,\dots,d.
	\]
	
	Defining $\eta^{(j)} := -(\mu_{j1},\dots,\mu_{jd})\in\mathbb{Z}^d$, 
	we obtain
	\begin{equation}\label{nDCs}
	|\eta^{(j)} + A(\boldsymbol{\omega})\xi^{(j)}| \le C_\varepsilon' e^{-\varepsilon| \xi^{(j)}|^{\frac{1}{s}}},\qquad\forall\,j\in\mathbb{N}.
	\end{equation}
	
	Since \eqref{nDCs} holds for every $\varepsilon>0$, $A(\boldsymbol{\omega})$ does not satisfy \eqref{DCss}.
	\end{proof}

\end{document}
